% Apêndice A: arquivos complementares para reprodução técnica do protótipo.
% O ambiente apendicesenv gera automaticamente "APÊNDICE A" no título e no sumário.
\begin{apendicesenv}
\chapter{Arquivos complementares para reprodução do protótipo}
\label{ap:arquivos_complementares}

Com a finalidade de disponibilizar os recursos utilizados no desenvolvimento do sistema e favorecer sua reprodução técnica, os arquivos complementares deste trabalho foram organizados em um repositório público na plataforma GitHub.

O repositório reúne o código-fonte do processamento de eventos MIDI em Python, o firmware do microcontrolador ESP32, os procedimentos experimentais implementados, a documentação das conexões elétricas, a relação de componentes e os arquivos tridimensionais utilizados na fabricação do invólucro. Também estão disponíveis arquivos CSV com resultados consolidados dos ensaios realizados.

No repositório, o código embarcado está em \path{firmware_esp32/}, os programas de controle e aquisição MIDI em \path{software_pc/}, a documentação elétrica e os modelos tridimensionais em \path{hardware/}, e os resultados consolidados em \path{dados_experimentais/}. As duas estratégias de mapeamento entre \textit{velocity} e PWM estão preservadas em programas separados, permitindo distinguir a implementação principal do procedimento de caracterização experimental.

\noindent\textbf{Repositório:} tcc-interface-musical-multissensorial

\noindent\textbf{Endereço eletrônico:}\par
\begingroup\sloppy\noindent\url{https://github.com/anaguiar777/tcc-interface-musical-multissensorial}\par\endgroup

\noindent\textbf{Acesso em:} 2 nov. 2026.

As orientações de instalação, configuração e execução dos programas encontram-se no arquivo \texttt{README.md} do repositório. A reprodução do protótipo exige a utilização de componentes compatíveis, a configuração da rede de comunicação e a verificação das conexões e tensões elétricas antes do acionamento.

Os materiais disponibilizados têm finalidade de documentação e reprodução técnica. Os dados experimentais consolidados não substituem integralmente os registros brutos das medições, e a reprodução do sistema não constitui, por si só, validação de sua eficácia perceptiva para pessoas surdas ou com deficiência auditiva.

\end{apendicesenv}
